% Auto-generated by python/build_tables.py
\begin{table*}[!htbp]
\centering
\begin{threeparttable}
\caption{Model constructs and managerial interpretation.}
\label{tab:constructs}
\small
\begin{tabular}{p{0.23\textwidth}p{0.23\textwidth}p{0.47\textwidth}}
\toprule
\textbf{Construct} & \textbf{Formal representation} & \textbf{Managerial interpretation} \\
\midrule
Relational confidence & $w_{ab}(t)\in[0,1]$ & Reliability of a recurring working relation between two agents. \\
Successful interaction & $x_{ab}(t)=1$ & An interaction that reinforces the working relation by improving understanding, alignment, or problem solving. \\
Failed interaction & $x_{ab}(t)=0$ & An interaction that weakens the working relation through misunderstanding, friction, or misalignment. \\
Boundary readiness & $R_B(t)$ & Fraction of cross-boundary ties that have reached sufficient relational confidence. \\
First-passage time & $T=\inf\{t:R_B(t)\geq q\}$ & Time required for the collaboration boundary to become relationally ready. \\
Translation advantage & $\pi_{BS}>\pi_{out}$ & Higher probability that boundary-spanning interactions succeed relative to ordinary cross-boundary interactions. \\
Boundary-spanner load & $k/(2b)$ & Approximate amount of cross-boundary work carried by each boundary spanner. \\
\bottomrule
\end{tabular}
\begin{tablenotes}
\footnotesize
\item Notes: The model does not represent full interpersonal trust or realised R\&D performance. It represents relational confidence at the tie level and relational coordination readiness at the boundary level.
\end{tablenotes}
\end{threeparttable}
\end{table*}
